\section{Research program, milestones, and falsification criteria}

The theory is implemented in stages. Each stage has a success criterion and a
falsification criterion so that numerical difficulty cannot be hidden by adding
more corrective layers.

\subsection{Stage A: continuous geometry kernel}

Deliver arbitrary-pose superelliptic spiral coils with finite superelliptic
cross-sections, superquadric/CSG packages, relative-pose scene graphs, robust
Bishop frames, collision/clearance tests, and high-order geometry quadrature.

\textbf{Success:} rigid-transform invariance and geometry derivatives agree with
independent finite-difference/analytic checks over broad poses and aspect
ratios.

\textbf{Falsification:} the representation develops discontinuities at shape
exponent changes, near-zero curvature, or large rotations.

\subsection{Stage B: homogeneous-medium conductor integral teacher}

Implement finite-cross-section current--charge/PEEC physics for one and multiple
coils in an unbounded homogeneous medium.

\textbf{Success:} impedance, current distribution, skin/proximity loss, and
mutual coupling converge under modal/quadrature enrichment and agree with an
independent converged reference on a locked benchmark set.

\textbf{Falsification:} accuracy requires degrees of freedom scaling like a full
three-dimensional background volume, or self terms require ad hoc
fine-minus-coarse corrections.

\subsection{Stage C: packages and general media}

Add SIE for piecewise homogeneous passive bodies, layered Green functions where
appropriate, and localized VIE/adaptive-node corrections for continuous
material contrast.

\textbf{Success:} homogeneous regions add interface rather than volume
complexity; local contrast complexity scales with contrast support.

\textbf{Falsification:} a remote homogeneous background change forces a global
volume remesh or world-grid state.

\subsection{Stage D: passive electrical macromodel}

Build narrowband and broadband structured port models and couple external
circuits.

\textbf{Success:} reciprocity/passivity are exact to numerical roundoff after
decoding and broadband time simulations remain stable without post-hoc pole
repair.

\textbf{Falsification:} small frequency-domain errors routinely create unstable
or active time-domain models.

\subsection{Stage E: continuous loss operator}

Construct continuous PSD conductor/dielectric loss factors and low-rank source
channels.

\textbf{Success:} non-negative local heat and integrated power closure hold for
arbitrary complex port currents; thermal observables converge with modest loss
rank.

\textbf{Falsification:} accuracy depends on a fixed world-cell ordering or
requires large projection corrections to restore positivity/power.

\subsection{Stage F: thermal transfer model}

Implement Green/BIE thermal teachers, stable reduced realizations, continuous
temperature decoders, and localized inhomogeneous corrections.

\textbf{Success:} transient and steady temperatures meet the declared error
budget across moving packages/media without global thermal remeshing.

\textbf{Falsification:} every new pose requires a large global matrix
factorization comparable to a full PDE solve.

\subsection{Stage G: neural correction and certified inference}

Train the scene operator, structured EM/thermal corrections, uncertainty
calibration, and predict-then-correct residual path.

\textbf{Success:} FAST inference is suitable for repeated optimization/control,
while CERTIFIED mode reaches reference accuracy with a small bounded correction
budget on most in-domain scenes.

\textbf{Falsification:} certification usually converges no faster than solving
the integral teacher from an uninformed initial guess.

\subsection{Stage H: full electrothermal deployment}

Couple temperature-dependent materials, electrical circuits, broadband or
carrier-envelope dynamics, and active learning.

\textbf{Success:} locked release suites cover pose, shape, package, material,
frequency, excitation, and time-evolution regimes with traceable error budgets.

\subsection{Primary research questions}

The following questions must be answered experimentally rather than assumed:
\begin{enumerate}
\item What cross-section modal family reaches skin/proximity accuracy with the
      smallest stable rank over the target frequency range?
\item Which hybrid conductor/SIE formulation gives the best conditioning for
      high conductivity contrast and closely spaced package interfaces?
\item How many continuous loss channels are required for thermal observables,
      and should they be optimized jointly with the thermal realization?
\item Which $\SE(3)$ scene architecture provides the best accuracy/latency trade
      for variable object count and strongly separated length scales?
\item Can a learned preconditioner reduce CERTIFIED corrections without
      weakening residual-based acceptance?
\item Which thermal Green/BIE reduction remains most robust under
      temperature-dependent material properties?
\end{enumerate}

The project should prefer falsifying a weak formulation early over preserving
compatibility with an expensive legacy implementation.
